% English version of paper/ja_v2/02_関連研究.md (2026-11-17).
\section{Related work}\label{sec:related}

\subsection{Semi-supervised domain adaptation}
Methods usable in semi-supervised domain adaptation, in which the target has a few labels, have been proposed for each information source: learning representations that do not reveal the domain (domain-adversarial neural networks, DANN) \cite{ganin2016dann} and aligning feature covariances (CORAL) \cite{sun2016coral} for unlabelled targets, entropy minimisation and pseudo-labels from semi-supervised learning \cite{grandvalet2004,lee2013pseudolabel}, and minimax entropy (MME) for semi-supervised domain adaptation \cite{saito2019mme}. Adding source data can also lower target performance (negative transfer) \cite{wang2019negative}. The method studies above compare methods that combine information sources and do not measure, in one setting, which part of the classifier error each source corrects.

\subsection{Cross-corpus MER with few labels}\label{sec:rel-cc}
For audio emotion classifiers that degrade on another corpus, adversarial learning that aligns representations with unlabelled target utterances \cite{abdelwahab2018dann} and learning that pulls together representations learned per corpus \cite{gideon2021addog} have been proposed. Ryumina et al.\ \cite{ryumina2025} evaluated the four corpora of our leave-one-corpus-out setting (IEMOCAP, MOSEI, MELD and AFEW; we use only the improvised part of IEMOCAP) with four emotions without target labels and showed that training on several corpora raised UAR by 1.7--6.6\pt. Lian et al.\ \cite{lian2026merbench} built cross-corpus benchmarks with shallow classifiers on frozen features and reported that multimodal models exceeded unimodal ones by small differences across corpora. For targets with a few labels, Feng and Chaspari \cite{feng2023} proposed metric learning for audio features, Li et al.\ \cite{li2021acted} adaptation of audio and video between acted and improvised speech, and He et al.\ \cite{he2025m4ser} a method using audio and speech-recognition transcripts.

\subsection{Distillation from other modalities}\label{sec:rel-distill}
Albanie et al.\ \cite{albanie2018} attached facial-expression predictions to unlabelled videos and trained an audio emotion classifier on them, and Zhang et al.\ \cite{zhang2021combining} combined this scheme with semi-supervised learning. Transferring representations or predictions of one modality to a model of another through paired unlabelled data is also used outside emotion recognition \cite{gupta2016,jaritz2020xmuda}. Training on the predictions of a more accurate model of another modality has been reported not to always improve the trained model \cite{xue2023mfh}.

\subsection{Late fusion and reliability weights}\label{sec:rel-fusion}
Zadeh et al.\ \cite{zadeh2017tfn} criticised late fusion, which trains unimodal classifiers separately and combines their predictions, for not learning interactions between modalities because it assumes that simple weighted averaging is a proper fusion approach. Predictions can be combined by product or sum rules \cite{kittler1998}. Busso et al.\ \cite{busso2004} reported that modality strengths differ by emotion, Kaya et al.\ \cite{kaya2017} chose model-by-emotion weights on AFEW by random search on validation data, M3ER \cite{mittal2020m3er} learned a multiplicative fusion layer that strengthens reliable modality cues per utterance, and Mar\'in et al.\ \cite{marin2026classwise} used per-utterance, per-emotion modality weights output by a network trained with all labels within one corpus.

\subsection{Position of this study}\label{sec:position}
% ver3: rewritten after reviews/v2_ja/prose_intro (Sect. 3) and citations_intro (16--18).
This study examines distillation in a setting that belongs to both the cross-corpus MER of Sect.~\ref{sec:rel-cc} and the distillation from other modalities of Sect.~\ref{sec:rel-distill}: cross-corpus MER with only a few target labels per emotion. Unlike Albanie et al., who used only facial predictions, our pseudo-labels are built from three classifiers including text. To the best of our knowledge, no prior study has compared, for the audio classifier under the same conditions, distillation on pseudo-labels including other modalities with 11 domain adaptation, semi-supervised and test-time adaptation methods that do not use such pseudo-labels (Sect.~\ref{sec:compmethods}; versions moved to linear classifiers on frozen features) in cross-corpus MER with audio, visual and text and only a few target labels per emotion. Three comparisons are close. Zhang et al.\ (2021) \cite{zhang2021combining} compared pseudo-labels including facial predictions with same-modality semi-supervised learning for speech emotion recognition, but used all training labels of the evaluated corpus (IEMOCAP) and did not compare domain adaptation methods. Jaritz et al.\ \cite{jaritz2023xmossda} compared cross-modal learning with uni-modal unsupervised domain adaptation methods moved to a semi-supervised setting, but for 3D semantic segmentation. Zhang et al.\ (2024) \cite{zhang2024amanda} proposed a multimodal domain adaptation method that uses no target labels (Amanda) between IEMOCAP, MELD, MOSEI and MSP-IMPROV and compared it with DANN and other domain adaptation methods, but not with distillation on pseudo-labels from other modalities.

Multimodal models trained on few samples can fall below the best unimodal model \cite{huang2021}. In RQ2 we use late fusion (Sect.~\ref{sec:rel-fusion}) to combine predictions into pseudo-labels. Unlike M3ER and Mar\'in et al., which learn per-utterance weights, the proposal uses one weight per classifier and emotion. To our knowledge, splitting a reliability table into classifier strengths measured on the K target labels and per-emotion profiles measured on the source corpora has not been reported.
